% TOP_PROBLEM: {"id": 3389, "title": "Counterexamples to the Baillie-PSW primality test", "category_id": 1, "category": {"id": 1, "name": "number_theory", "display_name": "Number Theory", "description": "Properties of integers, prime numbers, Diophantine equations.", "slug": "number-theory", "order_index": 1, "created_at": "2026-07-31T15:26:25.670Z"}, "status": "open"}
% FIRST_POSTED: null
% ENTRY_KIND: statement_only
% Imported from: batch04.tex; UnsolvedMath ID 3389
\documentclass[11pt]{article}
\usepackage[margin=25mm]{geometry}
\usepackage{fontspec}
\setmainfont{Latin Modern Roman}
\usepackage{amsmath,amssymb,mathtools}
\usepackage{unicode-math}
\setmathfont{Latin Modern Math}
\usepackage{xurl,hyperref}
\hypersetup{colorlinks=true,urlcolor=blue}
\setcounter{secnumdepth}{0}
\setlength{\parindent}{0pt}
\setlength{\parskip}{5pt}
\setlength{\emergencystretch}{3em}
\newcommand{\sref}[2]{\hyperlink{source:#1:#2}{[S#2]}}
\newcommand{\cataloguescope}[1]{\paragraph{Recorded target.}#1}
\begin{document}
% UnsolvedMath ID 3389; source catalogue code OPG-511
\setcounter{section}{3388}
\section{Counterexamples to the Baillie-PSW primality test}\label{3389}
{\small\sffamily UnsolvedMath ID 3389 · Number Theory}
\subsection{Definitions and mathematical statement}

\paragraph{Shared notation.}$\mathbb N=\{1,2,\ldots\}$, and $\mathbb Z,\mathbb Q,\mathbb R,\mathbb C$ denote the integers, rationals, reals and complex numbers. Any other conventions are specified locally.


For odd $n>1$, write $n-1=2^su$ with $u$ odd. The strong base-two test passes if $2^u\equiv1\pmod n$ or $2^{2^ju}\equiv-1\pmod n$ for some $0\le j<s$. For the Lucas test, choose $D$ successively from $5,-7,9,-11,\ldots$ until the Jacobi symbol $(D/n)=-1$; reject an encountered nontrivial common divisor of $D$ and $n$, and reject perfect squares. Put $P=1,Q=(1-D)/4$ (equivalently use Method A*, changing $P,Q$ to $5,5$ when $Q=-1$). Define $U_0=0,U_1=1$, $V_0=2,V_1=P$, and
\[
 U_{j+2}=PU_{j+1}-QU_j,\qquad V_{j+2}=PV_{j+1}-QV_j.
\]
Writing $n+1=2^tv$ with $v$ odd, the strong Lucas test passes when $U_v\equiv0\pmod n$ or $V_{2^jv}\equiv0\pmod n$ for some $0\le j<t$. The Baillie--PSW test combines these strong tests with the preceding rejections; even $n>2$ are rejected and $2$ is prime. Here $(D/n)$ is the Jacobi symbol, computed multiplicatively from the Legendre symbols at the prime factors of odd $n$. Finally, $F_0=0,F_1=1,F_{j+2}=F_{j+1}+F_j$ defines the Fibonacci numbers. \sref{3389}{3}
\paragraph{Statement recorded in the supplied source.}
Resolve both source questions:
\begin{enumerate}
\item Does some composite integer pass the Baillie--PSW test just specified? Exhibit one, or prove that none exists.
\item Does some composite $n\equiv3$ or $7\pmod{10}$ satisfy simultaneously
\[
 n\mid 2^{n-1}-1,\qquad n\mid F_{n+1}?
\]
Exhibit one, or prove nonexistence.
\end{enumerate}
The second question is stated separately, rather than identified with the strong-test condition. \sref{3389}{2}
\subsection{Short English statement}
Find a composite passing the specified combined primality test, or prove its universal correctness; separately settle the specified base-two/Fibonacci divisibility question.
\subsection{Sources}
\begin{description}
\item[\hypertarget{source:3389:1}{S1}] ULAM.AI and the UnsolvedMath Contributors, UnsolvedMath: A Curated Collection of Open Mathematics Problems, record 3389 (OPG-511). CC BY 4.0. \url{https://huggingface.co/datasets/ulamai/UnsolvedMath}.
\item[\hypertarget{source:3389:2}{S2}] Open Problem Garden, Counterexamples to the Baillie-PSW primality test, OPG-511. Original problem statement, full mathematical discussion, bibliography and comments. \url{https://www.openproblemgarden.org/op/counterexamples_to_the_baillie_psw_primality_test}.
\item[\hypertarget{source:3389:3}{S3}] Robert Baillie, Andrew Fiori and Samuel S. Wagstaff, Jr., Strengthening the Baillie--PSW primality test, arXiv:2006.14425v2 (2021). Sections 2.2--2.4 and 3, original strong test and equivalence of Method A and A*. \url{https://arxiv.org/abs/2006.14425v2}.
\end{description}
\label{end:3389}
\end{document}
